\section{Đạo hàm và đánh giá đa thức [Cơ bản]}
\subsection{Phân tích \& Thiết kế (M0)}
\begin{itemize}
    \item \textbf{Input:}
    \begin{itemize}
        \item File \texttt{polynomials.csv}: \texttt{id,degree,a0..an}
        \item File \texttt{x\_values.csv}: mỗi dòng là một giá trị \texttt{x}
    \end{itemize}

    \item \textbf{Xử lý:}
    \begin{itemize}
        \item \textbf{Case biên:}
        \begin{itemize}
            \item \texttt{degree < 0} $\rightarrow$ invalid
            \item Thiếu hệ số $\rightarrow$ gán bằng 0
            \item \texttt{degree = 0} $\rightarrow$ đạo hàm bằng 0
            \item \texttt{x\_values.csv} rỗng $\rightarrow$ chỉ ghi header
        \end{itemize}

        \item \textbf{Module 1 (Struct dữ liệu):}
\begin{lstlisting}
struct Poly {
    string id;
    int deg;
    vector<double> a;
};
\end{lstlisting}

        \item \textbf{Module 2 (Đọc polynomials.csv):}
        \begin{itemize}
            \item Chỉ lưu đa thức hợp lệ
        \end{itemize}

        \item \textbf{Module 3 (Đọc x\_values.csv):}
        \begin{itemize}
            \item Đọc các giá trị x
        \end{itemize}

        \item \textbf{Module 4 (Tính P(x) bằng Horner):}
\begin{lstlisting}
res = a[n];
for i = n-1 downto 0:
    res = res*x + a[i];
\end{lstlisting}

        \item \textbf{Module 5 (Tạo đa thức đạo hàm):}
        \[
        P(x)=a_0+a_1x+a_2x^2+\cdots+a_nx^n
        \]
        \[
        P'(x)=a_1+2a_2x+\cdots+na_nx^{n-1}
        \]

        \item \textbf{Module 6 (Ghi file):}
        \begin{itemize}
            \item Ghi ra \texttt{deriv\_report.csv}
            \item Định dạng:
            \texttt{id,x,P(x),dP(x)}
        \end{itemize}
    \end{itemize}

    \item \textbf{Output:}
\begin{lstlisting}
deriv_report.csv
id,x,P(x),dP(x)
...
\end{lstlisting}
\end{itemize}

\subsection{Mã nguồn C++11 (Tách module)}
\begin{lstlisting}
#include <iostream>
#include <fstream>
#include <sstream>
#include <vector>
#include <string>
#include <cstdlib>
#include <iomanip>

using namespace std;

struct Poly {
    string id;
    int deg;
    vector<double> a;
};

bool readPolysCSV(const string& path,
                  vector<Poly>& arr) {
    ifstream fin(path);
    if (!fin) return false;

    string line;

    while (getline(fin, line)) {
        if (line.empty()) continue;

        stringstream ss(line);
        vector<string> token;
        string x;

        while (getline(ss, x, ',')) {
            token.push_back(x);
        }

        if (token.size() < 2) continue;

        Poly p;
        p.id = token[0];
        p.deg = atoi(token[1].c_str());

        if (p.deg < 0) continue;

        p.a.assign(p.deg + 1, 0);

        for (int i = 0; i <= p.deg; ++i) {
            int pos = i + 2;

            if (pos < (int)token.size())
                p.a[i] = atof(token[pos].c_str());
            else
                p.a[i] = 0;
        }

        arr.push_back(p);
    }

    return true;
}

bool readXValues(const string& path,
                 vector<double>& xs) {
    ifstream fin(path);
    if (!fin) return false;

    string line;

    while (getline(fin, line)) {
        if (line.empty()) continue;

        xs.push_back(atof(line.c_str()));
    }

    return true;
}

double evalHorner(vector<double> a,
                  double x) {
    int n = a.size() - 1;

    double res = a[n];

    for (int i = n - 1; i >= 0; --i)
        res = res * x + a[i];

    return res;
}

vector<double> derivative(Poly p) {
    vector<double> d;

    if (p.deg == 0) {
        d.push_back(0);
        return d;
    }

    d.assign(p.deg, 0);

    for (int i = 1; i <= p.deg; ++i)
        d[i - 1] = i * p.a[i];

    return d;
}

int main() {
    vector<Poly> polys;
    vector<double> xs;

    if (!readPolysCSV("polynomials.csv",
                      polys)) {
        cerr << "Loi mo polynomials.csv!\n";
        return 1;
    }

    if (!readXValues("x_values.csv",
                     xs)) {
        cerr << "Loi mo x_values.csv!\n";
        return 1;
    }

    ofstream fout("deriv_report.csv");

    if (!fout) {
        cerr << "Loi mo output!\n";
        return 1;
    }

    fout << "id,x,P(x),dP(x)\n";
    fout << fixed << setprecision(6);

    for (size_t i = 0; i < polys.size(); ++i) {
        vector<double> d =
            derivative(polys[i]);

        for (size_t j = 0; j < xs.size(); ++j) {
            double val =
                evalHorner(polys[i].a, xs[j]);

            double der =
                evalHorner(d, xs[j]);

            fout << polys[i].id << ","
                 << xs[j] << ","
                 << val << ","
                 << der << "\n";
        }
    }

    cout << "Done";

    return 0;
}
\end{lstlisting}

\subsection{Kế hoạch kiểm thử (Test Cases)}

\textbf{Test Case 1 (Thông thường)}

\textbf{Input (polynomials.csv):}
\begin{lstlisting}
P1,2,1,2,1
\end{lstlisting}

\textbf{Input (x\_values.csv):}
\begin{lstlisting}
2
\end{lstlisting}

\textbf{Expected Output (deriv\_report.csv):}
\begin{lstlisting}
id,x,P(x),dP(x)
P1,2.000000,9.000000,6.000000
\end{lstlisting}

\vspace{0.5cm}

\textbf{Test Case 2 (Bậc 0)}

\textbf{Input (polynomials.csv):}
\begin{lstlisting}
A,0,5
\end{lstlisting}

\textbf{Input (x\_values.csv):}
\begin{lstlisting}
3
\end{lstlisting}

\textbf{Expected Output (deriv\_report.csv):}
\begin{lstlisting}
id,x,P(x),dP(x)
A,3.000000,5.000000,0.000000
\end{lstlisting}

\vspace{0.5cm}

\textbf{Test Case 3 (x rỗng)}

\textbf{Input (polynomials.csv):}
\begin{lstlisting}
B,1,1,1
\end{lstlisting}

\textbf{Input (x\_values.csv):}
\begin{lstlisting}

\end{lstlisting}

\textbf{Expected Output (deriv\_report.csv):}
\begin{lstlisting}
id,x,P(x),dP(x)
\end{lstlisting}